The first workers have received everything promised to them. What wounds them is the sight of someone else receiving it too. Their complaint brings a difficult question into the center of this Sunday: can we rejoice in a goodness that refuses to become our private advantage? The landowner's generosity exposes the difference between loving a gift and loving the superiority we think it gives us.

Isaiah calls the sinner to return to a God whose mercy exceeds the sinner's own thoughts. The psalm praises divine greatness together with divine nearness. Paul desires to be with Christ, yet recognizes that his continued labor will benefit others. The Gospel brings these questions to a point in a vineyard: work is real, the promise is kept, and the master's goodness is larger than the workers' calculation. The Mass's prayers ask for love to become obedience and for sacramental participation to bear fruit in conduct.

Three readings of the whole formulary develop that movement. Augustine's account of divine cultivation asks what God gives and what his gift makes possible in us. Chrysostom's account of the working day asks how those called early or late should respond now. Gregory the Great's account of the one vineyard asks how many generations belong to a people no generation owns. They agree that receiving God's goodness demands fruitful charity. They differ in their first identification of the vineyard, the meaning of its hours, and the explanation of the payment sequence. Each has something the others do not bring into equal focus.

\subsection{The Sunday from entrance to dismissal}
\begin{longtable}{@{}p{0.29\linewidth}p{0.66\linewidth}@{}}
\toprule
\textbf{Element} & \textbf{Appointment}\\\midrule
Entrance & \textit{Salus populi ego sum}; the Lord's saving help in distress.\\
Collect & \textit{Deus, qui sacrae legis}; the law of love and life with God.\\
First reading & Isaiah 55:6--9.\\
Responsorial Psalm & Psalm 145:2--3, 8--9, 17--18; response from 18a.\\
Second reading & Philippians 1:20c--24, 27a.\\
Gospel acclamation & Cf. Acts 16:14b; an adapted request for receptive hearing.\\
Gospel & Matthew 20:1--16a.\\
Prayer over the Offerings & \textit{Munera, quaesumus, Domine}.\\
Communion & Psalm 119:4--5 \textbf{or} John 10:14.\\
Prayer after Communion & \textit{Quos tuis, Domine}.\\\bottomrule
\end{longtable}
